\documentclass[11pt,a4paper, 
swedish, english %% Make sure to put the main language last!
]{article}
\pdfoutput=1

%% Andréas's custom package 
%% (Will work for most purposes, but is mainly focused on physics.)
\usepackage{../custom_as}

%% Figures can now be put in a folder: 
\graphicspath{ {figures/} %{some_folder_name/}
}

%% If you want to change the margins for just the captions
\usepackage[margin=10 pt]{caption}

%% To add todo-notes in the pdf
\usepackage[%disable  %%this will hide all notes
]{todonotes} 

%% Change the margin in the documents
\usepackage[
%            top    = 3cm,              %% top margin
%            bottom = 3cm,              %% bottom margin
%            left   = 2.5cm, right  = 2.5cm %% left and right margins
]{geometry}


%% If you want to change the formatting of the section headers
%\renewcommand{\thesection}{...}

\newcommand{\EM}{\text{EM}}
\newcommand{\E}{\text{E}}
%\newcommand{\GR}{\text{GR}}
%\newcommand{\SR}{\text{SR}}


%%%%%%%%%%%%%%%%%%%%%%%%%%%%%%%%%%%%%%%%%%%%%%%%%%%%%%%%%%%%%%%%%%%%%%
\begin{document}%% v v v v v v v v v v v v v v v v v v v v v v v v v v
%%%%%%%%%%%%%%%%%%%%%%%%%%%%%%%%%%%%%%%%%%%%%%%%%%%%%%%%%%%%%%%%%%%%%%


%% If you want to use an external file for the title page
%\renewcommand{\thefootnote}{\fnsymbol{footnote}}

%kortkommandon för mailaddresserna
\newcommand{\andsunds}{andsunds@student.chalmers.se}
\newcommand{\rigon}{rigon@student.chalmers.se}



\pagenumbering{roman} %%Romersk sidnumrering i början
\begin{titlepage}
\newgeometry{top=3cm, bottom=2cm}

\newcommand{\HRule}{\rule{\linewidth}{0.5mm}} % Defines a new command for the horizontal lines, change thickness here

\center % Center everything on the page
 
%------------------------------------------------------------------------------------
%	HEADING SECTIONS
%------------------------------------------------------------------------------------

\textsc{\huge Chalmers tekniska högskola}\\[1.5cm] % Name of university/college
\textsc{\Large Rapport, Experimentell fysik 2}\\[0.2cm] % Major heading such as course name
\textsc{\large Termodynamik -- Uppgift 3 }\\[0.5cm] % Minor heading such as course title

%------------------------------------------------------------------------------------
%	TITLE SECTION
%------------------------------------------------------------------------------------

\HRule \\[0.4cm]
{ \LARGE \bfseries 
Studier av kvicksilveratomens atomära emissionsspektra samt absorptionsspektra av laserfärgämnena Rhodamin B och Kumarin 307
}\\[0.4cm] % Title of  document
\HRule \\[1.5cm]
 
%------------------------------------------------------------------------------------
%	AUTHOR SECTION
%------------------------------------------------------------------------------------

\begin{minipage}{0.4\textwidth}
\begin{flushleft} \large
\emph{Författare:}\\
Andréas Sundström\footnotemark{} \\
Rigon Demisai\footnotemark{} 
\end{flushleft}
\end{minipage}
~
\begin{minipage}{0.4\textwidth}
\begin{flushright} \large
\emph{Labassistent:} \\
Martin Wersäll
\end{flushright}
\end{minipage}\\[3cm]

\setcounter{footnote}{0}
\stepcounter{footnote}
  \footnotetext{\href{mailto:\andsunds}{\texttt{\andsunds}}}
\stepcounter{footnote}
  \footnotetext{\href{mailto:\rigon}{\texttt{\rigon}}}



%------------------------------------------------------------------------------------
%	DATE SECTION
%------------------------------------------------------------------------------------
% Följer ISO-standarden för tidsintervall:
% https://en.wikipedia.org/wiki/ISO_8601#Time_intervals
% "Double hyphen" också ok istället för '/'. -- i LaTeX är dock lite på gränsen
{ \large
\begin{tabular}{rc}
    Laboration utförd: & 2015-12-11/15 \\[0.1cm]
    Rapport inlämnad: & \today
\end{tabular}\\[1cm]
}

%------------------------------------------------------------------------------------
%	LOGO SECTION
%------------------------------------------------------------------------------------

\includegraphics[height=5cm]{logo.pdf} % Include a department/university logo
 
%------------------------------------------------------------------------------------

\vfill % Fill the rest of the page with whitespace

\end{titlepage}
\restoregeometry


\setcounter{page}{2}%detta är ANDRA (2) sidan

\renewcommand{\abstractname}{Sammandrag}
\begin{abstract}
Den här rapporten beskriver en spektroskopisk studie av
kvicksilveratomens atomära emissionsspektrum samt en studie av
laserfärgämnena rhodamin~B och kumarin~307 och deras
absorptionsspektra.
Ur emisionspektrumet från Hg har vissa av atomens energinivåer kartlagts,
baserat på kvanmekaniska regler och jämförelse med tidigare data på
Hg.  Kvicksilvrets emissionsspektra
är taget i intervallet 350--1100\,nm, där detekterades totalt 23
emissionstoppar, varefter 13 övergångar kunde identifieras.
Absorptionsspektrumen för rhodamin~B och kumarin~307 visar breda
absorptionsband vilket är kännetecknande för flourescerande ämnen som
består av stora organiska molekyler. 
Mätningarna har utförts med en Spex~270M spektrometer och
datainsamlingen har gjorts i LabView.
\end{abstract}

\renewcommand{\abstractname}{Abstract}
\begin{abstract}
This report describes a spectroscopic study of the atomic emision
spectrum of mercury, and also a study of the laser dyes Rhodamine~D
and Coumarin~307 and their absorption spectrum.
Som of the energylevels of Hg have been idetified from the emission
spectrum, based on quantum mecanichal rules and comparison with erlier
data on Hg. The emission spectrum of mercury was taken in the interval
of 350--1100\,nm, there a total of 23 emission lines were detected,
from which 13 different atomic transitions could be identified. 
The absorption spectrum of Rhodamin~B and Coumarin~307  show wide
absorption bands which are characteristic of big organic molecules.
THe measurements were made with a Spex~270M spectrometer and the data
collection was done through LabVIEW.
\end{abstract}

\clearpage
\renewcommand{\contentsname}{Innehållsförteckning}
\tableofcontents

\clearpage
\pagenumbering{arabic}
\setcounter{page}{1}

\renewcommand{\thefootnote}{\arabic{footnote}}
\setcounter{footnote}{0}



%%%%%%%%%%%%%%%%%%%% vvv Internal title page vvv %%%%%%%%%%%%%%%%%%%%%
\title{Gravitation and Cosmology -- FFM 071
\\ {\Large Hand-in set 2} }
\author{Andréas Sundström}
\date{2018-03-06}

\maketitle

%%%%%%%%%%%%%%%%%%%% ^^^ Internal title page ^^^ %%%%%%%%%%%%%%%%%%%%%
%% If you want a list of all todos
%\todolist

\renewcommand{\thesubsection}{\arabic{section} (\roman{subsection})}
\setcounter{section}{8}

\section{3D anti-de Sitter}

The $\text{AdS}_3$ metric
\begin{equation}\label{eq9:metric1}
\rd{s}^2=-\rd{u}^2-\rd{v}^2+\rd{x}^2+\rd{y}^2
\end{equation}
under the constraint
\begin{equation}\label{eq9:constraint}
-u^2-v^2+x^2+y^2=-R^2,
\end{equation}
can be rewritten using the change of variables
\begin{equation}\label{eq9:var-change}
\begin{aligned}
&u=\sqrt{R^2+r^2}\cos(t/R)
\qc\qquad &x=r\cos\phi, \\
&v=\sqrt{R^2+r^2}\sin(t/R)
\qc\qquad &y=r\sin\phi.
\end{aligned}
\end{equation}
Firstly it is easy to confirm that \eqref{eq9:var-change} satisfy
\eqref{eq9:constraint} by noting that
\begin{equation}
u^2+v^2=R^2+r^2\qand
x^2+y^2=r^2.
\end{equation}
Then to calculate the metric, we need
\begin{equation}
\rd{u}=\pdv{u}{r}\rd{r}+\pdv{u}{t}\rd{t}+\pdv{u}{\phi}\rd\phi
=\frac{r\cos(t/R)\rd{r}}{\sqrt{R^2+r^2}}
-\frac{\sqrt{R^2+r^2}}{R}\sin(t/R)\rd{t}
\end{equation}
and similarly for
\begin{equation}
\rd{v}=\frac{r\sin(t/R)\rd{r}}{\sqrt{R^2+r^2}}
+\frac{\sqrt{R^2+r^2}}{R}\cos(t/R)\rd{t},
\end{equation}
which leads to
\begin{equation}
\rd{u}^2+\rd{v}^2=
\frac{r^2\rd{r}^2}{R^2+r^2}+\frac{R^2+r^2}{R^2}\rd{t}^2
\end{equation}
Then since we already recognize the polar coordinates for $x$ and $y$,
we can write down
\begin{equation}
\rd{x}^2+\rd{y}^2=\rd{r}^2+r^2\rd\phi^2.
\end{equation}
We can now write down the metric \eqref{eq9:metric1} in terms of the
new coordinates
\begin{equation}
\begin{aligned}
\rd{s}^2=&-\frac{R^2+r^2}{R^2}\rd{t}^2
+\qty[1-\frac{r^2}{R^2+r^2}]\rd{r}^2+r^2\rd\phi^2\\
=&-\frac{R^2+r^2}{R^2}\rd{t}^2
+\frac{R^2}{R^2+r^2}\rd{r}^2+r^2\rd\phi^2
\end{aligned}
\end{equation}

The proper distance from $r=0$ to $r=\infty$ can easily be calculated
as
\begin{equation}
S=\int_{r=0}^{r=\infty}\rd{s}=
\begin{Bmatrix}
\rd{t}=0\\
\rd{\phi}=0
\end{Bmatrix}
=\int_0^\infty\rd{r}
\overbrace{\frac{R}{\sqrt{R^2+r^2}}}^{\sim R/r\text{ for }r\gg R}
 =\infty.
\end{equation}
Meanwhile the time it takes for a light signal to traverses the same
path is
\begin{equation}
T=\int_{r=0}^{r=\infty}\rd{t}
=\int_0^\infty\rd{r}\frac{R^2}{R^2+r}
=\frac{\pi R}{2},
\end{equation}
where $\rd{t}$ have been calculated using
\begin{equation}
0=\rd{\tau^2}=\frac{R^2+r^2}{R^2}\rd{t}^2
-\frac{R^2}{R^2+r^2}\rd{r}^2
\end{equation}
for light.



\section{Maximally symmetric metric?}
We have a metric given by
\begin{equation}\label{eq10:ds}
\rd{s}^2=-\rd{t}^2+\rd{x}^2+\rd{y}^2+\rd{z}^2
-4\cosh(\frac{x}{2})
\qty[\cosh(\frac{x}{2})(\rd{t}+\rd{x})-\sinh(\frac{x}{2})\rd{y}]\rd{x},
\end{equation}
which in matrix form gives
\begin{equation}
g_{\mu\nu}=
\begin{bmatrix}
-1 & -(1+\cosh x) &0 &0\\
-(1+\cosh x) &-(1+2\cosh x) &+\sinh x&0\\
0&+\sinh x &1&0\\
0&0&0&1
\end{bmatrix}.
\end{equation}
Here we have used the hyperbolic identities
\begin{equation}
2\cosh^2\qty(\frac{x}{2})\equiv 1+\cosh x
\qand
2\cosh(\frac{x}{2})\sinh(\frac{x}{2})
\equiv \sinh x.
\end{equation}

To see if this is maximally symmetric, we calculate the Riemann
tensor. To do that we need the affine connection
\begin{equation}
\Gamma^\rho_{\mu\nu}=\frac{1}{2}g^{\rho\sigma}
\qty(\pd_{\mu}g_{\nu\sigma}+\pd_{\nu}g_{\mu\sigma}-\pd_{\sigma}g_{\mu\nu}).
\end{equation}
However we inside the parenthesis only the elements with $\mu\nu=xx$
survive. To prove this we take $\mu\nu=x\alpha$, where $\alpha\neq x$,
this gives
\begin{equation}
\pd_{x}g_{\alpha\sigma}+\cancelto{0}{\pd_{\alpha}}g_{x\sigma}-\pd_{\sigma}g_{x\alpha}
=\pd_{x}g_{\alpha\sigma}-\pd_{\sigma}g_{x\alpha}.
\end{equation}
The only way for either of the two remaining terms to survive is if
$\sigma=x$, but in that case the two terms are identical and cancel
each other. Then for $\mu\nu=\alpha\beta$, where $\alpha,\beta\neq x$,
we get
\vspace{-1em}
\begin{equation}
\cancelto{0}{\pd_{\alpha}}g_{\beta\sigma}
+\cancelto{0}{\pd_{\beta}}g_{\alpha\sigma}
-\cancel{\pd_{\sigma}g_{\alpha\beta}}=0.
\end{equation}
The last term cancel because the only non-constant elements in
$g_{\mu\nu}$ are on either $\mu=x$ or $\nu=x$, which is not the case
for $\mu\nu=\alpha\beta$. Hence the only non-vanishing elements of the
affine connection are $\Gamma^\rho_{xx}$. This then means that the
Riemann tensor,
\begin{equation}
{R^\rho}_{\sigma\mu\nu}=
\pd_{\mu}\Gamma^\rho_{\sigma\nu}-\pd_{\nu}\Gamma^\rho_{\sigma\mu}
+\Gamma^\rho_{\mu\tau}\Gamma^\tau_{\sigma\nu}
-\Gamma^\rho_{\nu\tau}\Gamma^\tau_{\sigma\mu},
\end{equation}
actually vanishes. To see this more clearly, we note that the
anti-symmetry in $\mu\nu$ means that $\mu\neq\nu$, but due to the
restriction to $\Gamma^\rho_{xx}$ we must have one of $\mu$ or $\nu$
to be $x$. Without loss of generality, we choose $\mu=x$ and
$\nu=\alpha\neq x$, and we get
\begin{equation}
{R^\rho}_{\sigma x\alpha}=
\pd_{x}\cancelto{0}{\Gamma^\rho_{\sigma\alpha}}
-\cancelto{0}{\pd_{\alpha}}\Gamma^\rho_{\sigma x}
+\Gamma^\rho_{x\tau}\cancelto{0}{\Gamma^\tau_{\sigma\alpha}}
-\cancelto{0}{\Gamma^\rho_{\alpha\tau}}\Gamma^\tau_{\sigma x}=0.
\end{equation}
This is in other words a flat space-time, and is therefore maximally
symmetric.
% Note to self: See Weinberg p. 380:
% "Thus any flat metric admits $N(N+1)/2$ independent Killing vectors,
% and is therefore maximally symmetric."


We will therefore now try to get \eqref{eq10:ds} into the familiar
form of the Minkowski metric,
\begin{equation}\label{eq10:Mink1}
\rd{s'}^2=-\rd{t'}^2+\rd{x'}^2+\rd{y'}^2+\rd{z'}^2.
\end{equation}
To do this we introduce the new variables $t'=t'(t,x)$, $x'=x'(x)$,
$y'=y'(x,y)$, and $z'=z'(z)=z$. For $z'$ the choice is obvious, and
for the other three variables it is relatively easy to convince
yourself that that they should only depend on $x$ and their non-primed
counterpart. We can now write out the metric \eqref{eq10:Mink1}
(ignoring $z'$ and $z$ since we already know that transformation) 
\begin{equation}%\label{eq10:Mink2}
\begin{aligned}
\rd{s'}^2=&
-\qty(\pdv{t'}{t})^2\rd{t}^2-2\pdv{t'}{t}\pdv{t'}{x}\id{t}\id{x}
-\qty(\pdv{t'}{x})^2\rd{x}^2\\
&+\qty(\pdv{x'}{x})^2\rd{x}^2\\
&+\qty(\pdv{y'}{x})^2\rd{x}^2+2\pdv{y'}{x}\pdv{y'}{y}\id{x}\id{y}
+\qty(\pdv{y'}{y})^2\rd{y}^2.
\end{aligned}
\end{equation}
For this to equal the original metric \eqref{eq10:ds}, we must have
\begin{equation}\label{eq10:cond}
\begin{cases}
\qty(\pdv{t'}{t})^2=1\\
\qty(\pdv{y'}{y})^2=1\\
\qty(\pdv{x'}{x})^2-\qty(\pdv{t'}{x})^2+\qty(\pdv{y'}{x})^2
=-(1+2\cosh x)\\
\pdv{t'}{t}\pdv{t'}{x}=1+\cosh x\\
\pdv{y'}{x}\pdv{y'}{y}=\sinh x.
\end{cases}
\end{equation}
If we choose the positive solutions to the first two equations, we get
\begin{equation}
\begin{cases}
\qty(\pdv{t'}{t})=+1\\
\qty(\pdv{y'}{y})=+1\\
\end{cases}
\quad\Longrightarrow\quad
\begin{cases}
t'(t,x)=t+f(x)\\
y'(x,y)=y+g(x),
\end{cases}
\end{equation}
which means that the last two equations in \eqref{eq10:cond} become
\begin{equation}
\begin{cases}
\dv{f}{x}=1+\cosh x\\
\dv{g}{x}=\sinh x
\end{cases}
\quad\Longrightarrow\quad
\begin{cases}
t'(t,x)=t+x+\sinh x\\
y'(x,y)=y+\cosh x.
\end{cases}
\end{equation}
It is now straight forward to show that the middle equation in
\eqref{eq10:cond} becomes
\begin{equation}
\qty(\pdv{x'}{x})^2
\cancel{-(1+2\cosh x)}
\overbrace{-\cosh^2 x+\sinh^2x}^{=-1}
=\cancel{-(1+2\cosh x)},
\end{equation}
which means that we also here can choose $\pdv*{x'}{x}=+1$. Every
thing is self-consistent, meaning that the full transformation of
variables 
\begin{equation}
\begin{cases}
t'=t'(t,x)=t+x+\sinh x\\
x'=x'(x)=x\\
y'=y'(x,y)=y+\cosh x\\
z'=z'(z)=z
\end{cases}
\end{equation}
results in $\rd{s'}^2$ from \eqref{eq10:Mink1} and $\rd{s}^2$ from
\eqref{eq10:ds} being equal. With this transformation of variables
the maximal symmetry is obvious.



\section{The non-flat torus}
On the non-flat torus embedded in three dimensions 
\begin{equation}
\begin{cases}
x=r_2A(\phi_2)\cos\phi_1\\
y=r_2A(\phi_2)\sin\phi_1\\
z=r_2\sin\phi_2
\end{cases}
\qq{with}
A(\phi_2)=r_1+r_2\cos\phi_2,
\end{equation}
and the metric, given by 
\begin{equation}
\rd{s}^2=\rd{x}^2+\rd{y}^2+\rd{z}^2
=A^2(\phi_2)\rd\phi_1^2+{r_2}^2\rd\phi_2^2,
\end{equation}
is
\begin{equation}
g_{\mu\nu}=
\begin{bmatrix}
A^2(\phi_2)&0\\
0&{r_2}^2
\end{bmatrix}.
\end{equation}
Here $\phi_1$ denotes an angle running around the ``hole'' of the
torus, and $\phi_2$ runs around the ``cylinder'' making up the torus;
$r_1$ and $r_2$ are the corresponding radii.
From this it is a simple task to compute the Affine connection
\begin{equation}
\Gamma^\rho_{\mu\nu}=\frac{1}{2}g^{\rho\sigma}
\qty(\pd_{\mu}g_{\nu\sigma}+\pd_{\nu}g_{\mu\sigma}-\pd_{\sigma}g_{\mu\nu}),
\end{equation}
whose only non-zero elements are 
\begin{equation}\label{eq11:affine}
\begin{aligned}
\Gamma^1_{12}=\Gamma^1_{21}
=&\frac{1}{2}g^{11}\pd_2g_{11}
=-\frac{r_2\sin\phi_2}{A(\phi_2)}\\
\Gamma^2_{11}=&-\frac{1}{2}g^{22}\pd_2g_{11}
=\frac{A(\phi_2)\sin\phi_2}{r_2}.
\end{aligned}
\end{equation}
Here we have used $\pd_2A^2(\phi_2)=-2A(\phi_2){r_2}\sin\phi_2$. Then
the Riemann tensor,
\begin{equation}
{R^\rho}_{\sigma\mu\nu}=
\pd_{\mu}\Gamma^\rho_{\sigma\nu}-\pd_{\nu}\Gamma^\rho_{\sigma\mu}
+\Gamma^\rho_{\mu\tau}\Gamma^\tau_{\sigma\nu}
-\Gamma^\rho_{\nu\tau}\Gamma^\tau_{\sigma\mu},
\end{equation}
is also seen to only have these few non-zero elements
\begin{equation}
\begin{aligned}
{R^1}_{212}=-{R^1}_{221}=&-\pd\Gamma^1_{21}-\Gamma^1_{21}\Gamma^1_{21}
=\frac{r_2}{A(\phi_2)}\cos\phi_2,\\
{R^2}_{112}=-{R^1}_{121}=&-\pd\Gamma^2_{11}+\Gamma^2_{11}\Gamma^1_{21}
=-\frac{A(\phi_2)}{r_2}\cos\phi_2.
\end{aligned}
\end{equation}
This gives the Ricci tensor
\begin{equation}
R_{\mu\nu}={R^\rho}_{\mu\rho\nu}={R^1}_{\mu1\nu}+{R^2}_{\mu2\nu}
=\begin{bmatrix}
\frac{A(\phi_2)}{r_2}\cos\phi_2&0\\
0&\frac{r_2}{A(\phi_2)}\cos\phi_2
\end{bmatrix},
\end{equation}
and the Ricci scalar
\begin{equation}
R=g^{\mu\nu}R_{\mu\nu}=g^{11}R_{11}+g^{22}R_{22}
=\frac{2\cos\phi_2}{r_2A(\phi_2)}
=\frac{2\cos\phi_2}{r_2(r_1+r_2\cos\phi_2)}.
\end{equation}
We now also note that $R_{\mu\nu}=g_{\mu\nu}R/2$ as is the case for
maximally symmetric spaces.

Now to the Killing vectors. We start with the Killing equations
\begin{equation}
0=\nabla_\mu\xi_\nu+\nabla_\nu\xi_{\mu}
=\pd_\mu\xi_\nu+\pd_\nu\xi_{\mu}-2\Gamma^\rho_{\mu\nu}\xi_\rho.
\end{equation}
This gives us three equations for $\mu\nu=11$, $12$, and $22$,
\begin{equation}\label{eq11:Killing0}
\begin{cases}
\pd_1\xi_1=\Gamma^\rho_{11}\xi_{\rho}
=\frac{A(\phi_2)\sin\phi_2}{r_2}\xi_{2}\\
\pd_1\xi_2+\pd_2\xi_1=2\Gamma^\rho_{12}\xi_{\rho}
=-2\frac{r_2\sin\phi_2}{A(\phi_2)}\xi_1\\
\pd_2\xi_2=\Gamma^\rho_{22}\xi_{\rho}=0.
\end{cases}
\end{equation}
Using the last equation we get
\begin{equation}
\xi_2=f(\phi_1),
\end{equation}
then substituting this into the first equation we get
\begin{equation}\label{eq11:xi1}
\xi_1=\frac{A(\phi_2)\sin\phi_2}{r_2}F(\phi_1)+g(\phi_2),
\end{equation}
where $\pdv*{F}{\phi_1}=f(\phi_1)$, and lastly the middle equation
becomes 
\begin{equation}
f'(\phi_1)+\qty(\frac{A(\phi_2)}{r_2}\cos\phi_2-\sin^2\phi_2)F(\phi_1)
+g'(\phi_2)=-2
\qty[\sin^2\phi_2F(\phi_1)+\frac{r_2\sin\phi_2}{A(\phi_2)}g(\phi_2)],
\end{equation}
where prime denotes differentiation with respect to the variable of the
respective functions. This can now be rewritten, using
$A(\phi_2)=r_1+r_2\cos\phi_2$, to 
\begin{equation}\label{eq11:Killing1}
f'(\phi_1)+g'(\phi_2)+\qty(1+\frac{r_1}{r_2}\cos\phi_2)F(\phi_1)
=-\frac{2r_2\sin\phi_2}{r_1+r_2\cos\phi_2}g(\phi_2).
\end{equation}
This equation must hold for any combination of $\phi_1$ and
$\phi_2$. Specifically if we fix $\phi_1$ at some value, we can
find a solution for $g(\phi_2)=\tilde{g}(\phi_2)$. However, this
solution must also be a solution for any other value of $\phi_1$, and
the only way this can be fulfilled is if we require that
$F(\phi_1)\equiv F_0$ is constant for all $\phi_1$. This also means
that $f(\phi_1)\equiv0$ and $f'(\phi_1)\equiv0$. We can now write
\eqref{eq11:Killing1} as
\begin{equation}
g'(\phi_2)+\frac{2r_2\sin\phi_2}{r_1+r_2\cos\phi_2}g(\phi_2)
=-\qty(1+\frac{r_1}{r_2}\cos\phi_2)F_0,
\end{equation}
which can be solved using the particular and homogeneous
solution. However to make the process a little more expedient we can
also use \textsl{Mathematica}, which gives the solution
\begin{equation}
g(\phi_2)=C\qty(r_1{+}r_2\cos\phi_2)^2
-F_0\frac{(r_1{+}r_2\cos\phi_2)\sin\phi_2}{r_2}
=CA^2(\phi_2)-F_0\frac{A(\phi_2)\sin\phi_2}{r_2}.
\end{equation}
The final Killing vectors now become
\begin{equation}
\xi_1=CA^2(\phi_2)\qc
\xi_2=0.
\end{equation}
Note that the second term in $g$ is canceled in \eqref{eq11:xi1}, and
the seeming second Killing vector does not exist. I.e. there is only
one Killing vector, and it corresponds to a translation in $\phi_1$;
in the embedded geometry that is a rotation around the ``hole'' of 
the torus, which obviously is the only symmetry of a torus.


\section{Wormholes}
Here we have the metric
\begin{equation}\label{eq12:ds}
\rd{s}^2=-\rd{t}^2+\rd{r}^2
+\qty(r^2+a^2)\qty(\rd\theta^2+\sin^2\theta\id\phi^2),
\end{equation}
giving 
\begin{equation}
g_{\mu\nu}=\text{Diag}
\begin{bmatrix}
-1,& 1,& (r^2+a^2),& (r^2+a^2)\sin^2
\end{bmatrix}.
\end{equation}
We now want to find the geodesic equations
\begin{equation}
\dv[2]{x^\mu}{\tau}+\Gamma^\mu_{\rho\nu}\dv{x^\rho}{\tau}\dv{x^\nu}{\tau}=0. 
\end{equation}
To do that we need the affine connection which, since $g_{\mu\nu}$ is
diagonal and only depends on $r$ and $\theta$, is relatively straight
forward to compute. Its only non-zero elements are
\begin{equation}\label{eq12:Gamma}
\begin{aligned}
&\Gamma^r_{\theta\theta}=-r\qc
\Gamma^r_{\phi\phi}=-r\sin^2\theta\\
&\Gamma^{\theta}_{r\theta}=\Gamma^{\theta}_{\theta r}
=\frac{r}{r^2+a^2}\qc
\Gamma^{\theta}_{\phi\phi}=-\sin\theta\,\cos\theta\\
&\Gamma^{\phi}_{r\phi}=\Gamma^{\phi}_{\phi r}=\frac{r}{r^2+a^2}\qc
\Gamma^\phi_{\theta\phi}=\Gamma^\phi_{\phi\theta}=\cot\theta,
\end{aligned}
\end{equation}
which gives the geodesic equations
\begin{equation}\label{eq12:geodesic}
\begin{cases}
\dv[2]{t}{\tau}=0\\
\dv[2]{r}{\tau}-r\qty(\dv{\theta}{\tau})^2
-r\sin^2\theta\,\qty(\dv{\phi}{\tau})^2=0\\
\dv[2]{\theta}{\tau}+\frac{2r}{r^2+a^2}\dv{r}{\tau}\dv{\theta}{\tau}
-\sin\theta\cos\theta\qty(\dv{\phi}{\tau})^2=0\\
\dv[2]{\phi}{\tau}+\frac{2r}{r^2+a^2}\dv{r}{\tau}\dv{\phi}{\tau}
+\cot\theta\dv{\phi}{\tau}\dv{\theta}{\tau}=0.
\end{cases}
\end{equation}
The last two equations can be rewritten, by multiplying with
$(r^2+a^2)$ (and also with $\sin^2\theta$ for the last equation), to
\begin{equation}
\begin{cases}
\dv{\tau}\qty[\qty(r^2+a^2)\dv{\theta}{\tau}]
=\qty(r^2+a^2)\sin\theta\cos\theta\qty(\dv{\phi}{\tau})^2\\[0.5ex]
\dv{\tau}\qty[\qty(r^2+a^2)\sin^2\theta\dv{\phi}{\tau}]=0.
\end{cases}
\end{equation}
The last equation represents a conservation law similar to the $z$
component of the angular momentum
\begin{equation}\label{eq12:J}
\qty(r^2+a^2)\sin^2\theta\dv{\phi}{\tau}=J,
\end{equation}
where $J$ is a constant of motion. We also have
\begin{equation}\label{eq12:k}
\dv{t}{\tau}=k
\end{equation}
from the first equation of \eqref{eq12:geodesic}, where $k$ is
constant. Now for completeness sake, we note that the metric is
asymptotically Minkowski, which means that for a massive particle
\begin{equation}
k=\eval{\dv{t}{\tau}}_{r\to\infty}=
\eval{\dv{t}{\tau}}_{\text{Minkowski}}=\gamma=\qty(1-v^2)^{-1/2},
\end{equation}
where $v$ is the asymptotic velocity of the particle as $r\to\infty$.

Now to see if there are time-like geodesics traversing the wormhole,
we can restrict our search to the $\theta=\pi/2$ plane since the
metric is spherically symmetric\footnotemark{}. In this plane, we have
\footnotetext{See \eqref{eq12:ds-sphere} and the discussion
  immediately preceding it.} 
\begin{equation}
\rd\tau^2=\rd{t}^2-\rd{r}^2-\qty(r^2+a^2)\rd{\phi}^2,
\end{equation}
which when dividing through with $\rd\tau^2$ and substituting in
\eqref{eq12:J} and \eqref{eq12:k} becomes
\begin{equation}
+1=k^{2}-\qty(\dv{r}{\tau})^2-\frac{J^2}{r^2+a^2}.
\end{equation}
The $+1$ on the LHS is for time-like geodesics, as we are looking
for. We may now further simplify our search by restricting ourselves
to ``radial'' geodesics, i.e. with no angular momentum $J=0$. This
yields
\begin{equation}
\dv{r}{t}=\dv{\tau}{t}\dv{r}{\tau}=\pm\frac{1}{k}\sqrt{k^2-1}
=\pm v,
\end{equation}
which means that
\begin{equation}
r(t)=\pm vt+r_0;
\end{equation}
a supringly simple result.
This geodesic clearly has the capability to traverse the wormhole, in
the sense that it can go from arbitrarily large negative values of $r$
to arbitrarily large positive values of $r$. 


Using \eqref{eq12:Gamma} we can easily calculate the Riemann tensor
whose non-zero elements are
\begin{equation}
\begin{aligned}
&{R^r}_{\theta r\theta}=-{R^r}_{\theta\theta r}
=-\frac{a^2}{r^2+a^2} \qc
{R^r}_{\phi r\phi}=-{R^r}_{\phi\phi r}
=-\frac{a^2\sin^2\theta}{R^2+a^2}\\
&{R^\theta}_{r\theta r}=-{R^\theta}_{rr\theta}
=-\frac{a^2}{(r^2+a^2)^2}\qc
{R^\theta}_{\phi\theta\phi}=-{R^\theta}_{\phi\phi\theta}
=\frac{a^2\sin^2\theta}{r^2+a^2},\\
&{R^\phi}_{r\phi r}=-{R^\phi}_{rr\phi}
=-\frac{a^2}{(r^2+a^2)^2}\qc
{R^\phi}_{\theta\phi\theta}=-{R^\phi}_{\theta\theta\phi}
=\frac{a^2}{r^2+a^2}.
\end{aligned}
\end{equation}
From this we can compute the Ricci tensor and scalar. It is clear from
the Riemann tensor above that $R_{\mu\nu}={R^\rho}_{\mu\rho\nu}$ must
be diagonal and is given by
\begin{equation}
R_{\mu\nu}=\text{Diag}
\begin{bmatrix}
0,& -\frac{2a^2}{(r^2+a^2)^2},& 0, &0
\end{bmatrix},
\end{equation}
and then the curvature tensor trivially becomes
\begin{equation}
R=g^{\mu\nu}R_{\mu\nu}=-\frac{2a^2}{(r^2+a^2)^2}.
\end{equation}
We also see that $a\to0$ recovers a flat space, as we would expect
from the metric.

We now can write down the Einstein tensor as
\begin{equation}
G_{\mu\nu}=R_{\mu\nu}-\frac{R}{2}g_{\mu\nu}
=\frac{a^2}{(r^2+a^2)^2}\,\text{Diag}\!
\begin{bmatrix}
-1,& -1,&
(r^2+a^2),&
(r^2+a^2)\sin^2\theta
\end{bmatrix}.
\end{equation}
There are however problems with this Einstein tensor. According to
Einstein's equations 
\begin{equation}
G_{\mu\nu}=8\pi G T_{\mu\nu},
\end{equation}
and therefore since $G_{00}<0$ also $T_{00}<0$, but $T_{00}$
represents the energy density which can not be negative. 

To find the isometries of this metric, we begin by noting that this
metric has a maximally symmetric subspace. Namely a spherically
symmetric sub-spaces, which is seen by writing \eqref{eq12:ds} on the
form 
\begin{equation}\label{eq12:ds-sphere}
\rd{s}^2=\hat{g}_{ab}\rd{x^a}\rd{x^b}
+f(t,r)\tilde{g}_{ij}\rd{x^i}\rd{x^j},
\end{equation}
where $a,b\in\qty{t,x}$, $i,j\in\qty{\theta,\phi}$, and
$\tilde{g}_{ij}$ is the maximally symmetric metric of the
2-sphere. This means that we immediately have the isometries of the
2-sphere, i.e. rotations around any axis. Apart from these isometries,
we can refine our analysis of \eqref{eq12:ds-sphere} by noting that
$f$ actually is also independent of $t$, and that $\hat{g}_{ab}$ is
constant (i.e. independent of $t$ and $r$) and diagonal. This means
that we can also do a parity transformation of $t$ and $r$, as well as
translate $t$; these isometries are therefore
\begin{equation}
\begin{cases}
t\to t'=\pm t+t_0\\
r\to r'=\pm r.
\end{cases}
\end{equation}
We cannot translate or boost in space, since the workhole is located
at $r=0$ so any of those two transformations will obviously not be
isometric. We can therfore not go any further unless $a=0$, and that
would just result in the regular Minkowski metric which is maximally
symmetric. 





%%%%%%%%%%%%%%%%%%%%%%%%%%%%%%%%%%%%%%%%%%%%%%%%%%%%%%%%%%%%%%%%%%%%%%
\end{document}%% ^ ^ ^ ^ ^ ^ ^ ^ ^ ^ ^ ^ ^ ^ ^ ^ ^ ^ ^ ^ ^ ^ ^ ^ ^ ^ ^
%%%%%%%%%%%%%%%%%%%%%%%%%%%%%%%%%%%%%%%%%%%%%%%%%%%%%%%%%%%%%%%%%%%%%%



%%% Local Variables:
%%% mode: latex
%%% TeX-master: t
%%% End:

%  LocalWords:  contravariant stereographic covariant variational
%  LocalWords:  azimuthal affine parametrization Ricci Minkowski
%  LocalWords:  Mathematica isometries
